% =========================================================
%  Merged: main_overleaf_template.tex + results_chapter.tex
% =========================================================
\documentclass[12pt,a4paper]{article}

% ---- encoding & language --------------------------------
\usepackage[T1]{fontenc}
\usepackage[utf8]{inputenc}
\usepackage[icelandic]{babel}

% ---- page layout ----------------------------------------
\usepackage[margin=2.5cm]{geometry}
\usepackage{graphicx}
\usepackage{fancyhdr}
\usepackage{setspace}

% ---- header / footer ------------------------------------
\pagestyle{fancy}
\fancyhf{}
\lhead{\includegraphics[height=1.2cm]{myndir/logo_nnv}}
\renewcommand{\headrulewidth}{0pt}
\cfoot{\thepage}
\setlength{\headheight}{40pt}

% ---- bibliography ---------------------------------------
\usepackage[
  backend=biber,
  style=authoryear,
  sorting=nyt,
  maxcitenames=2,
  mincitenames=1,
  maxbibnames=99,
  autolang=other,
  babel=other
]{biblatex}

\addbibresource{references.bib}

\renewbibmacro{in:}{}

\DefineBibliographyStrings{icelandic}{
  and       = {og},
  andothers = {o.fl.},
  editor    = {ritstjóri},
  editors   = {ritstjórar},
  byeditor  = {ritstj.},
  urlseen   = {Sótt}
}

\renewcommand*{\multinamedelim}{\addcomma\space}
\renewcommand*{\finalnamedelim}{\space og\space}
\renewcommand*{\finalandcomma}{}

% ---- tables (from results chapter) ----------------------
\usepackage{booktabs}
\usepackage{longtable}
\usepackage{array}
\usepackage{multirow}
\usepackage{colortbl}
\usepackage{makecell}
\usepackage{xcolor}

% ---- figures --------------------------------------------
\usepackage{float}
\floatplacement{figure}{H}
\usepackage{caption}
\captionsetup{font=small}
\usepackage{wrapfig}

% ---- misc pandoc helpers --------------------------------
\usepackage{amsmath,amssymb}
\setlength{\emergencystretch}{3em}
\providecommand{\tightlist}{%
  \setlength{\itemsep}{0pt}\setlength{\parskip}{0pt}}

% ---- hyperref -------------------------------------------
\usepackage{hyperref}
\hypersetup{hidelinks}

% =========================================================
\begin{document}

% ---- title block ----------------------------------------
\vspace*{1cm}

\begin{center}
    {\Large \textbf{Minnisblað}}
\end{center}

\vspace{1cm}

\noindent
\begin{tabular}{@{}l l@{}}
\textbf{Dagsetning:} & 9. apríl 2026 \\
\textbf{Viðtakandi:} & Aníta Ósk Áskelsdóttir \\
\textbf{Sendandi:}   & Starri Heiðmarsson og Valtýr Sigurðsson \\
\textbf{Efni:}       & Vistgerðir, gagnagæði og vísbendingar um nýtingu þeirra til vöktunar \\
\end{tabular}

\vspace{1cm}

\noindent
Landvistgerðir voru afmarkaðar og þeim lýst á landsvísu í fjölriti NÍ frá
2016 \cite{ottosson2016} en áður höfðu farið fram rannsóknir og flokkun
vistgerða á hálendi sbr. skýrslur frá nokkrum svæðum á miðhálendinu
(\cite{einarsson2000} \& \cite{magnusson2002}). Þegar vöktun
náttúruverndarsvæða var hafin í samstarfi þáverandi Náttúrufræðistofnunar
Íslands og náttúrustofa þá var ákveðið að hefja vöktun á vistgerðum með
því að endurmæla úrval þeirra sniða sem áður höfðu verið mæld. Í þessu
minnisblaði er farið yfir framkvæmdina, gerð tilraun til að bera saman
mælingar og meta vöktunargildið auk þess að leggja endurbætur á
framkvæmdahliðinni.

\vspace{1cm}

\noindent
Þau snið sem metin voru á árabilinu 1999--2013 voru mismunandi nema hvað
þau voru öll 400 m$^2$ en ýmist 200 m eða 40 m að lengd. Munur liggur
einnig í því að lagðir voru út 8 smáreitir á 200 m sniðunum en einungis
sex á þeim sem eru 40 m að lengd. Staðsetning sniðanna voru skráð með gps
tæki, upphafs-, loka- og miðhnit í 200 m sniðunum en einungis upphafs- og
lokahnit í 40 m sniðunum. Það sem skiptir kannski mestu varðandi nákvæmni
vöktunar er að ekki er hægt að staðsetja sniðin nákvæmlega vegna
ónákvæmni gps hnitanna, reynt hefur verið að ráða bót á þessu til
framtíðar með því að setja tréhæla við upphaf og lok þeirra sniðan sem
hafa verið endurmæld.

\vspace{1cm}

\noindent
Tölfræðileg úrvinnsla var gerð í R (v4.4.2). Fyrir hvern reit var
tegundafjöldi (S) reiknaður út frá pöruðum gögnum fyrri og síðari
sýnatöku og breyting á tegundafjölda ($\Delta$S) reiknuð sem mismunur
milli heimsóknanna tveggja. Breyting á tegundasamsetningu (BC) í hverjum
reit var mæld með Bray--Curtis-fjarlægð milli tegundasamsetningar fyrri og
seinni sýnatöku. Wilcoxon signed-rank próf var notað til að kanna hvort
breyting á tegundafjölda væri marktækt frábrugðin núlli yfir reitina í
heild. Spearman-fylgni var reiknuð milli $\Delta$S og
Bray--Curtis-fjarlægðar annars vegar og þriggja umhverfisþátta hins vegar:
hæðar yfir sjávarmáli, jarðvegsdýptar og tíma milli sýnataka.

\vspace{1cm}

\noindent
\textbf{Fjölbreytnistuðlar.} Líffræðileg fjölbreytni er mikið til umræðu
en hingað til hefur einkum verið litið til tegundafjölbreytni þegar
vistgerðir eru annars vegar. Hér prófuðum við fleiri ólíka stuðla,
þ.e.\ „Shannon", „Simpson" og „evenness". Kom í ljós allmikill munur milli
mælinga en þar skipti árafjöldi milli mælinga miklu máli.

\begin{figure}[H]
    \centering
    \includegraphics[width=0.8\textwidth]{myndir/fjolbreytni.png}
    \caption{Myndatexti}
\end{figure}

\newpage

% =========================================================
%  NIÐURSTÖÐUR
% =========================================================

\section{Minnkun tegundafjölda}

Tegundafjöldi minnkaði á milli fyrri og seinni mælingar í 183 af 305
reitum (60\%), á meðan 102 reitir (33\%) sýndu aukningu og 20 héldust
stöðugir á milli athugana. Miðgildi breytingarinnar var $-$2 tegundir á
reit (meðaltal $-$2,2). Wilcoxon-próf staðfesti að minnkunin er marktækt
frábrugðin núlli (V~=~12.113,5, p~$<$~0,001), sem þýðir að hnignunin er
raunveruleg og kerfisbundin, ekki tilviljanakennd.

\begin{figure}[H]
\centering
\includegraphics[width=0.9\textwidth]{myndir/richness-density-1}
\caption{Dreifing tegundafjölda við fyrri (blátt) og seinni (rautt)
mælingu. Lóðrétta línan sýnir miðgildi fyrri mælinganna.}
\end{figure}

% ---------------------------------------------------------
\section{Endurnýjun tegunda (Bray--Curtis)}

Endurnýjun tegunda var mæld með Bray--Curtis-fjarlægð milli fyrri og
seinni mælinga í hverjum reit. Miðgildi Bray--Curtis-gildis var 0,609
(meðaltal 0,585, n~=~163 reitir), sem þýðir að um helmingur þyngdaðrar
tegundasamsetningar breyttist á milli mælinga að meðaltali. Fylgni milli
endurnýjunar tegunda og minnkunar tegundafjölda var marktæk (Spearman
$\rho~=~-0,323$, p~$<$~0,001), en fremur veik --- sem bendir til þess að
tvær ólíkar tegundir breytinga séu í gangi samtímis: nettótap tegunda
annars vegar og skipti á tegundasamsetningu hins vegar.

\begin{figure}[H]
\centering
\includegraphics[width=0.9\textwidth]{myndir/bc-gap-1}
\caption{Breyting á tegundasamsetningu (Bray--Curtis) sem fall af tíma
milli mælinga. Rauð brotalína~= línuleg aðlögun (95\,\% CI); blá
lína~= loess-sléttun.}
\end{figure}

Breyting á tegundasamsetningu jókst marktækt með auknum tíma milli mælinga
(hallatala~=~$+$0,025 BC-einingar/ár, $R^2$~=~0,541, p~$<$~0,001).
Loess-kúrfan sýnir að breytingin er hvað hröðust á fyrstu 15 árunum og
hægist síðan eftir um 20 ár.

% ---------------------------------------------------------
\section{Umhverfisþættir sem spá fyrir um breytingu}

Þrír umhverfisþættir voru metnir sem mögulegar skýribreytur: hæð yfir
sjávarmáli, jarðvegsdýpi og fjöldi ára milli mælinga. Hrá Spearman-fylgni
sýnir að hæð og tími milli mælinga fylgjast mjög vel að ($\rho \approx
0,64$--$0,70$ fyrir BC), en þessir þættir eru ekki óháðir: reitir á
hálendinu höfðu almennt lengra bil milli mælinga. Hlutfylgni, þar sem
áhrif hinna þriggja breyta eru hafðar fastar, sýnir hvernig þættirnir
hegða sér óháð hverjum öðrum (tafla~\ref{tab:pred}).

\begin{table}[H]
\centering
\caption{Spearman-fylgni og hlutfylgni þriggja umhverfisþátta við
$\Delta$S og BC. Hlutfylgni eru reiknuð með öllum þremur þáttum í
líkaninu samtímis. Feitletraðar dálkar sýna hlutfylgni.}
\label{tab:pred}
\begin{tabular}{lrrrrrr}
\toprule
Þáttur &
  $\Delta$S $\rho$ &
  \textbf{$\Delta$S $\rho_p$} &
  $\Delta$S p &
  BC $\rho$ &
  \textbf{BC $\rho_p$} &
  BC p \\
\midrule
Hæð yfir sjávarmáli      & $-$0,343 & \textbf{$-$0,134} & $<$0,001 &  0,643 & \textbf{0,170} & $<$0,001 \\
Jarðvegsdýpi              & $-$0,110 & \textbf{$-$0,188} & 0,056    & $-$0,017 & \textbf{0,131} & 0,768 \\
Bil milli mælinga (ár)   & $-$0,349 & \textbf{$-$0,106} & $<$0,001 &  0,700 & \textbf{0,378} & $<$0,001 \\
\bottomrule
\end{tabular}
\end{table}

Þrjár niðurstöður skera sig úr. \textbf{Tími milli mælinga} hefur mest
áhrif á florábreytingu (BC hlutfylgni $\rho_p~=~0,378$, p~$<$~0,001), sem
staðfestir að gróðurbreytingin er enn í gangi og safnast upp með tíma.
\textbf{Hæð yfir sjávarmáli} er marktæk óháð tíma milli mælinga
($\Delta$S $\rho_p~=~-0,134$; BC $\rho_p~=~0,170$), sem bendir til þess
að hálendisreitir séu viðkvæmari. \textbf{Jarðvegsdýpi} kemur fram sem
sterkasta sjálfstæða skýribreytan fyrir \emph{tegundatap} ($\Delta$S
$\rho_p~=~-0,188$, p~$<$~0,05) --- grunnjörð tapar fleiri tegundum en
djúp jörð.

\begin{figure}[H]
\centering
\includegraphics[width=0.85\textwidth]{myndir/soil-depth-plot-1}
\caption{Tegundatap sem fall af jarðvegsdýpi. Rauð lína~= loess-sléttun;
blá brotalína~= línuleg aðlögun.}
\end{figure}

% ---------------------------------------------------------
\section{Jarðvegsbleyta}

\begin{figure}[H]
\centering
\includegraphics[width=\textwidth]{myndir/moisture-maps-1}
\caption{Florábreyting (BC) eftir jarðvegsbleytuflokkum. Grænt~= lítil
breyting; appelsínugult~= miðlungs; rautt~= mikil breyting.}
\end{figure}

\begin{figure}[H]
\centering
\includegraphics[width=\textwidth]{myndir/moisture-maps-2}
\caption{Breyting á tegundafjölda ($\Delta$S) eftir
jarðvegsbleytuflokkum. Dökkrautt~= mikið tap; appelsínugult~= lítið tap;
grátt~= stöðugt; blátt~= aukning.}
\end{figure}

\begin{table}[H]
\centering
\caption{Florábreyting og tegundatap eftir jarðvegsbleytuflokkum.
Kruskal--Wallis: BC $\chi^2$~=~13,12, p~=~0,004;
$\Delta$S $\chi^2$~=~11,06, p~=~0,011.}
\label{tab:moist}
\begin{tabular}{lrrr}
\toprule
Jarðvegsbleyta & n reitir & Miðgildi BC & Miðgildi $\Delta$S \\
\midrule
Þurrlendi  & 111 & 0,678 & $-$2 \\
Mýri       &  77 & 0,609 & $-$4 \\
Deiglendi  &  70 & 0,526 & $-$1 \\
Flói       &  45 & 0,441 & $-$1 \\
\bottomrule
\end{tabular}
\end{table}

Þurrlendi sýndi hæstu florábreytingu (BC miðgildi~=~0,678) en fremur
lítið tegundatap ($\Delta$S miðgildi~=~$-$2), sem bendir til þess að
tegundasamsetningin sé að breytast --- sumar tegundir hverfa en aðrar koma
í staðinn --- án þess að heildarfjöldi tegunda minnki verulega. Mýrar
sýndu hins vegar mest tegundatap ($\Delta$S~=~$-$4) án tilsvarandi hækkunar
á Bray--Curtis-fjarlægðinni, sem bendir til þess að tegundir séu að hverfa
án þess að nýjar taki við. Þetta eru tveir grundvallarólíkir
breytingarferlar sem skipta máli fyrir túlkun niðurstaðnanna.

% ---------------------------------------------------------
\section{Vistgerðir}

\begin{figure}[H]
\centering
\includegraphics[width=\textwidth]{myndir/habitat-maps-1}
\caption{Florábreyting (BC) eftir helstu vistgerðum (n~$\geq$~8 reitir).}
\end{figure}

\begin{figure}[H]
\centering
\includegraphics[width=\textwidth]{myndir/habitat-maps-2}
\caption{Breyting á tegundafjölda ($\Delta$S) eftir helstu vistgerðum
(n~$\geq$~8 reitir).}
\end{figure}

\begin{table}[H]
\centering
\caption{Florábreyting og tegundatap eftir vistgerð (n~$\geq$~8).
Kruskal--Wallis: BC $\chi^2$~=~81,05, p~$<$~0,001;
$\Delta$S $\chi^2$~=~53,60, p~$<$~0,001.}
\label{tab:hab}
\begin{tabular}{lrrr}
\toprule
Vistgerð & n reitir & Miðgildi BC & Miðgildi $\Delta$S \\
\midrule
Lyngmóavist á hálendi  &  8 & 0,802 & $-$8,0 \\
Rústamýravist          & 31 & 0,798 & $-$7,0 \\
Rekjuvist              & 17 & 0,682 & $-$4,0 \\
Starungsmýravist       & 34 & 0,543 & $-$3,5 \\
Sjávarfitjungsvist     & 27 & 0,379 &  $+$1,0 \\
Mosahraunavist         &  8 & 0,182 &  $+$7,0 \\
\bottomrule
\end{tabular}
\end{table}

Þrjár vistgerðir skera sig úr. \textbf{Rústamýravist} (n~=~31) sýndi
bæði mikla florábreytingu (BC~=~0,798) og mikið tegundatap
($\Delta$S~=~$-$7) og er þannig sú vistgerð sem er hvað verst sett
samkvæmt báðum mælikvörðum. \textbf{Lyngmóavist á hálendi} sýndi
svipaðar niðurstöður (BC~=~0,802, $\Delta$S~=~$-$8) og endurspeglar
líklega hæðaráhrifin sem komu fram í hlutfylgnigreiningu. Á hinn bóginn
stendur \textbf{Mosahraunavist} upp úr sem eina vistgerðin sem sýnir bæði
lága florábreytingu (BC~=~0,182) og aukningu á tegundafjölda
($\Delta$S~=~$+$7) og kann að gegna hlutverki sem friðland í breytu
landslagi.

% =========================================================
\printbibliography[title={Heimildir}, heading=subbib]

\end{document}
